\documentclass[12pt]{article}
\usepackage[T1]{fontenc}
\usepackage[utf8]{inputenc}
\usepackage{lmodern}
\usepackage{textcomp}
\usepackage{enumitem}
\usepackage{geometry}
\usepackage{graphicx}
\usepackage{amsmath, amssymb}
\geometry{margin=1in}
\begin{document}
\begin{center}
Biology - Undergraduate Level Assessment 17 \
Total Marks: 40 \
Answer all questions.
\end{center}

\section*{Multiple Choice (10 marks)}
\begin{enumerate}[label=\arabic*.]
\item A DNA molecule has 180 base pairs and is 20 percent adenine. How many cytosine nucleotides are present inthis moleculeof DNA?
\begin{enumerate}[label=(\alph*)]
    \item 108 nucleotides
    \item 120 nucleotides
    \item 150 nucleotides
    \item 144 nucleotides
    \item 60 nucleotides
\end{enumerate}
\item What hormones are involved with the changes that occur in a pubescent female and male?
\begin{enumerate}[label=(\alph*)]
    \item Cortisol, aldosterone, and adrenal androgens
    \item Insulin, glucagon, and somatostatin
    \item Cholecystokinin, ghrelin, and leptin
    \item Vasopressin, erythropoietin, and renin
    \item Epinephrine, norepinephrine, and dopamine
\end{enumerate}
\item About 0.5 percent of the American population have epilepsy, the falling sickness. Discuses the genetics of this disease.
\begin{enumerate}[label=(\alph*)]
    \item Epilepsy is the result of a chromosomal aneuploidy, similar to Down syndrome.
    \item Epilepsy is always inherited in an X-linked recessive pattern.
    \item Epilepsy is a dominant trait.
    \item Epilepsy is caused by an autosomal dominant trait.
    \item Epilepsy is due to an autosomal recessive trait.
\end{enumerate}
\item In the formation of the earliest cells, which of the following components most likely arose first?
\begin{enumerate}[label=(\alph*)]
    \item Lysosome
    \item Chloroplast
    \item Ribosome
    \item Golgi apparatus
    \item Plasma membrane
\end{enumerate}
\item Which of the following characteristics would allow you to distinguish a prokaryotic cell from an animal cell?
\begin{enumerate}[label=(\alph*)]
    \item Centrioles
    \item Chloroplasts
    \item Lysosomes
    \item Cell wall
    \item Mitochondria
\end{enumerate}
\end{enumerate}

\section*{True / False (10 marks)}
\textit{Answer True or False}
\begin{enumerate}[label=\arabic*.]
\setcounter{enumi}{5}
\item True or False: Polygenic traits are characteristics that are influenced by multiple genes, resulting in a wide range of phenotypes.
\item True or False: Cycads possess the ability to photosynthesize at night, which is a common trait among many plant species.
\item True or False: Complete equilibrium in a gene pool can always be achieved under ideal conditions with no external influences.
\item True or False: In chick embryos, the yolk sac is the extraembryonic membrane that provides nourishment to the developing fetus.
\item True or False: Distantly related organisms can have different rates of mutation due to varying evolutionary pressures and environmental factors.
\end{enumerate}

\section*{Long Form Response (20 marks)}
\textit{Answer each question with a clear and complete explanation.}
\begin{enumerate}[label=\arabic*.]
\setcounter{enumi}{10}
\item Discuss the number of chromatids present during meiosis in an organism with a diploid number of 8, specifically in the tetrad stage, late telophase of the first meiotic division, and metaphase of the second meiotic division.
\item Discuss how temperature, precipitation, soil chemistry, and decomposer types influence the characteristics and diversity of organisms within an ecosystem.
\end{enumerate}

\vfill
\noindent\textit{End of Paper}
\end{document}